\uuid{uZU3}
\exo7id{4103}
\titre{exo7 4103}
\auteur{quercia}
\organisation{exo7}
\datecreate{2010-03-11}
\isIndication{false}
\isCorrection{true}
\chapitre{Equation différentielle}
\sousChapitre{Equations différentielles linéaires}
\cpge{mpsi-mp2i.primitives-equations-differentielles.b, pcsi.techniques-calcul-primitives.b, ptsi.techniques-calcul-primitives.b, mpsi-mp2i.primitives-equations-differentielles.a, pcsi.techniques-calcul-primitives.a, ptsi.techniques-calcul-primitives.a, bcpst1.equations-differentielles-lineaires.a, mp-mpi.equations-differentielles-lineaires.a}
\module{Analyse}
\niveau{L2}
\difficulte{}

\contenu{
\question{
On pose $f(x) =  \int_{t=0}^{+\infty} e^{-t}e^{itx} \frac {d t}{\sqrt t}$.
Former une équation différentielle satisfaite par $f$. En déduire $f$.
}
\reponse{
$f'(x) = \frac {x+i}{2(x^2+1)} f(x)
  \Rightarrow  f(x) = \sqrt\pi (x^2+1)^{-1/4}\exp\left(\frac i2\arctan x\right)$.
}
}
